\section{Representative internal \(\ipra\) targets}
\label{app:ipra-target}

This appendix spells out the distinction between the finite exact kernel and a
statement about represented real coordinates.

\subsection{A single exact rational checker theorem}

Let \(w\in\N\) encode a tuple of signed rational values for all symbols in
(4.16).  Define primitive-recursive characteristic functions:
\begin{itemize}
\item \(\mathsf{CoreCert}(w)=1\) when all denominator codes are positive and
all equations in (4.16) hold after cross multiplication;
\item \(\mathsf{CoreConsequences}(w)=1\) when the cross-multiplied side and
angle cosine laws, sine law, unit-angle equations, five-part identities and
Cagnoli equations all hold.
\end{itemize}
Both functions use only bounded tuple decoding and natural arithmetic.  The
first exact target is
\begin{equation}
 \ipra\vdash
 \forall w\,
 \bigl(\mathsf{CoreCert}(w)=1\to
       \mathsf{CoreConsequences}(w)=1\bigr).                \tag{B.1}
\end{equation}
A verified polynomial normalizer can produce the derivation of (B.1).  No real
numbers or limits occur here.

For the tetrahedron, an analogous checker
\(\mathsf{TetCert}\) verifies four coordinate face vectors, their closure,
positive squared norms and the normalized dihedral dot equations.  Its target
is
\begin{equation}
 \ipra\vdash
 \forall w\,
 \bigl(\mathsf{TetCert}(w)=1\to
       \mathsf{TetCosine}(w)=1\bigr).                       \tag{B.2}
\end{equation}

\subsection{A fixed-name finite-precision theorem}

Now fix, in the metatheory, particular primitive-recursive programs for the
nine coordinates of \(u,v,w\), their moduli and bounds.  The programs are not
object-language variables.  Compile particular PR symbols for all output names
and let \(\mathsf{ApproxHyp}(n,\vec q)\) assert that the rational tuple
\(\vec q\) consists of the finitely many input approximations requested by the
precision schedule for output precision \(n\).  Let
\(\mathsf{CosineError}(n,\vec q)\) be the natural-number inequality obtained
by clearing positive denominators in
\begin{equation}
 \abs{\cos a-
 (\cos b\cos c+\sin b\sin c\cos A)}\le2^{-n}.              \tag{B.3}
\end{equation}
The actual internal theorem has the form
\begin{equation}
 \ipra\vdash
 \forall n\,\forall\vec q\,
 \bigl(\mathsf{ApproxHyp}(n,\vec q)
 \to\mathsf{CosineError}(n,\vec q)\bigr).                  \tag{B.4}
\end{equation}
The proof at free \(n\) instantiates input regularity only at explicitly
computed indices, invokes the finite polynomial identity, and closes a rational
error estimate.  Its induction calls are on bounded loops and use
quantifier-free invariants.  Universal introduction of \(n\) comes last.

The sine law, half formulas and all remaining identities use the same shape.
A quotient conclusion additionally assumes a numerical apartness certificate;
a positive-root conclusion includes the bounded-square-root bracket and the
branch interval.

\subsection{Lean object-language shape}

Using the actual repository syntax, the intended endpoint is schematically
\begin{lstlisting}[language={}]
open LeanProofs.ConciseIPRA

def sphericalHypotheses : Fm 0 := ...
def sphericalConclusions : Fm 0 := ...

theorem sphericalAtlas :
  Prf PRAAx [] (Fm.imp sphericalHypotheses sphericalConclusions) := by
  -- concrete PR defining equations
  -- qf induction for bounded arithmetic and error loops
  -- internally sound polynomial certificates
  -- ordinary intuitionistic introduction/elimination rules
  ...
\end{lstlisting}
The ellipsis denotes work, not an assumed theorem.  The supplied
\code{IPRABridge.lean} defines the data interfaces and target signatures so
that the remaining implementation obligations are visible rather than hidden
behind host-level axioms.

\subsection{Semantic presentation theorem}

After the internal theorem is complete, the existing ProveIt
bi-interpretation can present it in the concise native language.  Separately,
a semantic theorem states that standard rapid dyadic names satisfying the
premises denote the ordinary geometric quantities and hence validate the
familiar real-number formulas.  The direction of justification is therefore
\[
 \text{finite arithmetic derivation}
 \longrightarrow
 \text{name equality at every precision}
 \longrightarrow
 \text{ordinary geometric presentation},
\]
not the reverse.
